\documentclass[11pt,letterpaper]{article}
\usepackage[margin=1in]{geometry}
\usepackage[T1]{fontenc}
\usepackage{lmodern}
\usepackage{microtype}
\usepackage{parskip}
\usepackage[hidelinks]{hyperref}
\pagestyle{empty}
\setlength{\parindent}{0pt}
\begin{document}
15 November 2026

Editors\\
\emph{Pattern Recognition Letters}

Dear Editors,

Please consider \emph{Global loss bounds and one-bit recovery for direction-normalized half-moons} as a research report in \emph{Pattern Recognition Letters}. The manuscript establishes a global classification-loss criterion for direction normalization and a minimal constructive repair for a continuous benchmark.

For absolutely continuous binary populations in feature dimension at least two, we prove that directional Bayes risk converges uniformly at distant origins to orthogonal-projection risk. This closes the unbounded origin search and characterizes a positive global loss floor. In the plane, connected regions favoring opposite decisions with mutual convex-hull intrusion force this floor; robust signed-density cores yield explicit numerical lower certificates. The half-moon application covers the entire proved disjoint displacement interval under arbitrary positive within-support densities and priors, together with an overlapping configuration.

At displacement 0.5, one common radius threshold about the midpoint exactly restores the class decision. Since direction alone loses decisions at every origin, one additional radial bit is minimal for this population. The result includes an explicit decoder, a complete geometric proof and the exact interval of thresholds for which its two angular rules preserve every class decision. That interval also yields a sharp guarantee for bounded radius-measurement error and determines the threshold with the largest symmetric tolerance. A fresh controlled neural experiment produced three errors across 500,000 test observations with the bit, compared with 54,362 errors for direction alone. The full empirical record comprises 500 candidate fits and 250 validation-selected evaluations across the original and added studies.

The paper connects statistical decision sufficiency, geometric normalization limits and constrained information retention. It credits established wedge geometry, projections at infinity and classification-aware quantization, and develops the uniform risk criterion, explicit all-origin certificate and exact radial repair. Complete proofs, certified calculations, frozen protocols, trained models, predictions and reproducible analyses accompany the manuscript.

Thank you for considering the manuscript.

Sincerely,\\[0.6em]
Meshaal Mouawad\\
Department of Electrical and Computer Engineering\\
Mississippi State University\\
Mississippi State, MS 39762, USA
\end{document}
